\chapter{系统测试}
\section{系统测试概述}
系统测试的目标是验证 Search Agent 在真实工程场景下的功能正确性、性能指标、可靠性与可用性，确保从 prompt 到检索/爬取、解析、清洗、入库与返回的端到端流水线在各种正常与异常条件下满足设计要求。本章测试策略采用“分层+闭环”方法：先通过单元测试保证最小功能单元的正确性，再通过功能测试验证各核心模块在集成环境下的行为，最后通过非功能测试评估系统的工程可用性。为了兼顾工程效率与现实覆盖，测试采用混合数据源：一方面使用合成与回放历史真实请求来进行自动化回归；另一方面构造标注测试集与人工审查场景来评估检索质量与解析准确性。所有测试组件都纳入 CI/CD 流水线，关键阶段需通过灰度发布并触发自动化回归验证。

测试环境分为单元测试环境、集成测试环境与性能/压力释放环境。单元与集成测试在 Kubernetes 测试命名空间中以 mock 外部依赖方式进行，以便快速反馈并保证可重复性；性能与压力测试在独立大规模集群中运行，能模拟真实并发、网络延迟与外部系统行为。测试过程强调可观测性：所有测试运行时需采集 Prometheus 指标、日志与链路追踪，并将结果自动汇总到测试报告系统。质量门控策略将基于预先定义的 SLI决定是否允许发布或回滚。下文按测试类型详细展开测试设计与实现细节。
\section{单元测试}
单元测试着重验证系统最小功能单元的正确性与边界行为，是保证系统稳定性的第一道防线。针对 Search Agent 的不同模块，单元测试覆盖点包括：Prompt 解析模块的输入预处理、模型输出后处理与规则融合逻辑；Search Agent 的路由决策器；实时适配器与爬取任务构造器的任务描述格式与幂等性校验；数据处理模块的解析规则、字段标准化与去重算法；索引客户端的批量写入拼接与错误重试逻辑；召回融合与重排模块的分数归一化、特征计算与模型推理调用契约等。单元测试采用 Mock 与 Stub 技术隔离外部依赖。
\begin{figure}[H]
\centering
\includegraphics[
width=1.\textwidth,keepaspectratio]{picc/5.pdf}
\caption{单元测试代码}
\label{fig:zhouqishixian}
\end{figure}

实现细节方面，所有模块需编写覆盖率高的测试用例，覆盖正常输入、边界输入、异常输入（非法字符、恶意注入）与错误路径（模型超时、后端返回 5xx、解析异常）。Prompt 解析的单元测试不仅断言关键词与 intent 的正确性，还需检验置信度阈值分支、fallback 路径与规则覆盖情况。重排模块的单元测试需确保特征计算的数值稳定、score normalization 在多个通道输入下的稳定性与去重逻辑的正确性。数据库层的单元测试采用内存化数据库或 Docker Compose 启动轻量 DB，验证事务、唯一约束与回滚逻辑。
\begin{table}[htb]
\centering
\caption{系统单元测试结果表}
\label{tab:unit-test-system}
{\renewcommand{\arraystretch}{1.2}
\begin{tabular}{p{4cm} p{7cm} p{2cm}}
\hline
功能模块 & 测试内容 & 通过率 (\%) \\
\hline
Prompt解析模块 & 关键词提取准确性与边界输入鲁棒性测试 & 95\% \\ \hline
MCP通信与超时处理 & 请求/响应超时、重试与错误回退场景测试 & 96\% \\ \hline
实时检索（在线） & 热点定时检索与检索结果实时推送正确性 & 93\% \\ \hline
实时爬取与解析 & 实时爬取、解析正确性及超时重试测试 & 90\% \\ \hline
离线检索与索引 & 批量抓取、快照持久化与索引构建正确性 & 92\% \\ \hline
索引增量更新 & 增量合并、一致性与回滚测试 & 94\% \\ \hline
召回融合与重排 & 多源合并、去重与重排序准确性测试 & 91\% \\ \hline
数据处理与清洗 & 初次清洗、二次清洗与质量评估规则测试 & 89\% \\ \hline
向量检索性能 & 向量查询召回率与响应时延性能测试 & 88\% \\ \hline
端到端集成测试 & 从Prompt解析到数据落库的全流程验证 & 90\% \\ \hline
\end{tabular}
}
\end{table}
测试自动化覆盖 CI：每次 PR 都触发单元测试，测试失败阻止合并；同时在模型或规则变更提交时触发特定解析单元测试套件，以避免解析误回归。单元测试结果与覆盖率阈值作为代码质量核查项，并定期对测试用例做漏测补充（基于错误回放与事故 RCA）。
\section{功能测试}
功能测试在集成层面验证模块间接口契约、数据流与端到端功能场景。功能测试环境尽量接近生产，使用真实或高拟真度的外部依赖仿真，并利用历史请求回放作为主要测试驱动。以下分模块详细描述功能测试策略与评估准则。
\subsection{Prompt 解析模块功能测试} 
Prompt 解析模块的功能测试旨在验证解析质量、稳定性与与下游契约的一致性。测试采用标注语料库与回放真实交互日志两条主线：标注语料覆盖常见场景（问答型、命令型、检索型、模糊描述）、多意图示例、噪声输入、以及触发爬取的样例。每个输入标注预期输出，作为黄金标准。测试时将黄金输入送入在线解析服务，校验输出字段的精确度、召回率与置信度分布，并对比历史模型/规则输出以评估迭代改进带来的改动。
\begin{table}[htb]
\centering
\caption{Prompt解析模块功能测试报告}
\label{tab:prompt-parsing-test-brief}
{\renewcommand{\arraystretch}{1.2}
\begin{tabular}{p{2.5cm} p{3cm} p{2cm} p{3.5cm} p{1cm}}
\hline
测试功能 & 前置条件 & 测试步骤 & 预期结果 & 结果 \\
\hline
关键词提取 & 模型可用 & 输入含实体的prompt & 返回关键词列表 & 通过 \\\hline
空输入处理 & 接口可访问 & 传入空的字符串 & 返回空或明确错误码 & 通过 \\\hline
噪声鲁棒性 & 支持长文本 & 输入长文并夹杂噪声 & 抽出核心关键词 & 通过 \\\hline
实时路径判定 & 触发词规则已载入 & 输入含“详情”等词 & 标记为realtime=true & 通过 \\\hline
MCP超时重试 & 可模拟延迟/失败 & 模拟首次调用超时 & 重试或降级 & 通过 \\\hline
敏感词过滤 & 敏感词表已加载 & 输入含敏感词的prompt & 敏感词被屏蔽或替换 & 通过 \\\hline
响应性能 & 并发测试环境 & 并发多请求测延迟 & 平均响应<200ms & 通过 \\\hline
\end{tabular}
}
\end{table}
回放日志测试用于覆盖长尾与真实业务分布。将历史 prompt（匿名与脱敏）回放到解析模块并记录解析后下游影响指标：解析触发爬取比率、解析导致的检索通道路由变化、以及后续检索质量。对解析引起的误触率做重点关注，因为误触会显著增加爬虫池负载与外部风险。测试过程中通过断言与阈值设定（如误触率 < 0.1\%）判断是否通过。

功能测试还需验证解析模块的回退与版本管理路径：包括当在线模型超时或返回低置信度时，规则回退到何种逻辑、与离线强模型的回流校正是否能正确更新规则与训练样本、以及解析输出的 trace 链路与审计记录是否完整。最后，压力下的功能稳定性测试检验缓存命中、批处理推理、以及降级策略是否按预期触发。
\subsection{实时检索与爬取模块功能测试}
实时检索与爬取功能测试覆盖从 Search Agent 发起检索、判定触发爬取、到爬虫执行、解析并将结果下发的完整闭环。测试场景包括正常流、必须爬取的场景、目标站点 JS 渲染场景、以及故障与合规降级场景。功能测试需要部署仿真第三方站点：静态页面、需要渲染的页面、含反爬机制的页面的样例，以验证爬虫策略的自动检测与处理。
\begin{table}[htb]
\centering
\caption{实时检索与爬取模块功能测试报告}
\label{tab:realtime-retrieval-test-brief}
{\renewcommand{\arraystretch}{1.2}
\begin{tabular}{p{2.5cm} p{3cm} p{2cm} p{3.5cm} p{1cm}}
\hline
测试功能 & 前置条件 & 测试步骤 & 预期结果 & 结果 \\
\hline
热点定时检索 & 调度服务已启动 & 触发定时任务 & 返回候选 & 通过 \\\hline
在线检索请求 & 在线后端可用 & 发起检索请求 & 返回候选列表 & 通过 \\\hline
检索结果流推送 & 后端支持流式推送 & 接收并转发检索流 & 推送到消费端 & 通过 \\\hline
关键词触发爬取 & 爬虫服务可用 & 对含触发词候选触发爬取 & 爬虫开始抓取 & 通过 \\\hline
爬取超时与重试 & 可模拟延迟/失败 & 模拟首次超时场景 & 按重试策略重试或上报失败 & 通过 \\\hline
实时解析与抽取 & 解析组件在线 & 爬取后执行解析流程 & 提取标题/正文等 & 通过 \\\hline
解析结果推送 & 数据处理端可用 & 发送解析结果并确认消费 & 消息入队并被消费 & 通过 \\\hline
\end{tabular}
}
\end{table}
重点验证点包括爬取任务从生成到完成的 lifecycle（去重、入队、执行、回调）、任务幂等性（重复提交不导致重复入库）、断点续传（对大页面或网络抖动后的恢复能力）、与超时与重试策略的正确触发。测试会模拟目标站点临时不可用，观察调度器是否按退避策略重试、队列堆积是否触发速率限制、以及 Search Agent 是否在短时间窗口内进行合适的降级（仅返回摘要或链接）。对于渲染场景，测试要验证渲染池的资源隔离、渲染超时的中止与后续解析失败标记。

同时，需要验证爬取输出进入数据处理的契约：解析字段是否按 schema 填充、parse\_confidence 的分布是否合理、异常解析是否被正确标记并进入人工复核流程。合规测试包括在爬取前核验 robots.txt、在被禁止场景下是否正确跳过并记录审计日志、以及在触及敏感信息时是否正确脱敏或触发复核。对于触发爬取的同一 URL 频繁请求场景，要验证去重与短期缓存逻辑能有效抑制冗余抓取。
\subsection{离线检索与索引模块功能测试}
离线检索与索引的功能测试重点在于大规模数据管线的正确性、索引一致性与增量更新流程。测试流程从模拟数据采集开始：通过导入样本批次到对象存储，触发 ETL 流程并验证清洗、去重、标签化及向量化的输出。测试要断言清洗规则对噪声 HTML、缺失字段与非中文文本的处理是否符合设计，验证去重算法在常见重复变体上的召回与误判率。
\begin{table}[htb]
\centering
\caption{离线检索与索引模块功能测试报告}
\label{tab:offline-retrieval-test-brief}
{\renewcommand{\arraystretch}{1.2}
\begin{tabular}{p{3cm} p{2cm} p{2cm} p{3cm} p{2cm}}
\hline
测试功能 & 前置条件 & 测试步骤 & 预期结果 & 测试结果 \\
\hline
批量抓取调度 & 调度与爬虫可用 & 触发定时批量抓取 & 生成并持久化快照 & 通过 \\\hline
快照获取 & 目标域数据存在 & 请求指定域快照 & 返回完整域快照列表 & 通过 \\\hline
倒排索引构建 & 文本分词器可用 & 批量文档建索引 & 文本检索能正确命中 & 通过 \\\hline
向量索引构建 & 向量化服务在线 & 批量生成向量并入库 & 向量检索返回相似结果 & 通过 \\\hline
增量索引更新 & 索引维护服务可用 & 提交增量变更集并合并 & 新增/更新数据生效可搜到 & 通过 \\\hline
离线文本检索 & 倒排索引可查询 & 执行关键词查询 & 返回相关文档并排序合理 & 通过 \\\hline
离线向量检索 & 向量索引可查询 & 使用query向量检索 & 返回TopK相似结果 & 通过 \\\hline
\end{tabular}
}
\end{table}
向量化测试涉及 embedding 模型版本一致性检验、向量索引构建 correctness、以及增量插入逻辑。倒排索引测试需要验证映射兼容性，并在索引切换时通过 alias 原子替换检查查询不中断。索引发布前需要进行回归查询测试：对一组标准查询集合计算 Recall@K、nDCG 与 baseline 比较，确保新索引在核心指标上不降级。

此外，增量更新场景测试检验 delta 索引并行查询的正确性，以及合并窗口与合并后索引一致性。最后，测试索引回滚流程：在发现新索引性能问题时，验证别名回退、数据完整性校验以及重新建索引的可重复性与时间成本。
\subsection{召回融合与重排模块功能测试}
召回融合与重排的功能测试聚焦多通道候选合并、去重、特征计算、模型推理与最终排序质量。测试输入由多路候选源产生：实时倒排、离线倒排、向量召回与爬取解析结果，测试要验证归一化策略在合并后的稳健性，确保没有单一通道过度主导结果。去重测试检验两阶段去重在 doc\_id、URL 变体与语义近似重复上的表现，避免重要但类似的条目被误去除。
\begin{table}[htb]
\centering
\caption{召回融合与重排模块功能测试报告}
\label{tab:recall-rerank-test-brief}
{\renewcommand{\arraystretch}{1.2}
\begin{tabular}{p{3cm} p{2cm} p{2cm} p{3cm} p{2cm}}
\hline
测试功能 & 前置条件 & 测试步骤 & 预期结果 & 测试结果 \\
\hline
多源候选合并 & 多个候选源 & 提供多源列表并调用合并 & 所有候选被合并无丢失 & 通过 \\\hline
去重策略验证 & 存在重复候选 & 输入重复候选测试去重 & 重复被去除仅保留首项 & 通过 \\\hline
基线打分与排序 & 基线打分器可用 & 对候选执行基线打分 & 返回分数并按分降序排序 & 通过 \\\hline
模型重排效果 & 重排模型在线 & 将打分候选送入重排模型 & 返回重排序列表且顺序变化 & 通过 \\\hline
多样性过滤 & 存在单域刷量候选 & 执行多样性过滤算法 & 结果域分布更均衡 & 通过 \\\hline
阈值截断与TopK & 阈值与K配置好 & 执行截断与TopK截取 & 返回TopK且低分被过滤 & 通过 \\\hline
日志与指标上报 & 指标系统接入 & 运行全流程并检查指标 & 产出召回与最终数量指标 & 通过 \\\hline
\end{tabular}
}
\end{table}
重排功能测试需覆盖特征完整性、特征缺失下的容错性，以及排序模型推理稳定性。测试用离线标注集评价最终排序指标与人工标注结果的相关性，同时在多种负载与延迟限制下校验分层重排逻辑：当 cross-encoder 超时或不可用时系统是否能无缝降级到 LightGBM 并保持合理效果。还需验证排序解释性输出与溯源字段是否正确，以支持审计与人工校验。

重排上线前通过离线回放、A/B 测试与小流量灰度进行验证：离线回放用于快速筛查模型回归，A/B 实验与灰度检验用于在线测量 CTR、用户满意度及系统资源影响，同时统计显著性与长期趋势用于决策是否全量发布。
\section{非功能测试}
非功能测试关注系统的性能、稳定性、伸缩性、安全性与恢复能力。首先，性能测试在独立集群中执行：通过流量生成器模拟真实请求分布，测量关键 SLIs：P50/P95/P99 延迟、系统吞吐（QPS）、爬取并发数、消息队列 lag、索引入库吞吐与 CPU/内存/IOPS 等资源指标。压力测试逐步提升负载直至系统瓶颈，用于定位限流点与容量规划。性能优化点通常包括模型推理吞吐、重排深度控制、缓存命中率提升与批量写入优化。

稳定性测试包括长时运行的 soak 测试，在接近生产工作负载下观察内存泄露、慢查询累积、segment 合并导致的延迟陡增等问题。还要进行混沌测试：注入网络延迟、模拟 bin search 或数据库短暂不可用、强制重启爬虫池节点、制造磁盘IO峰值等，验证系统的降级策略与自动恢复能力。恢复能力测试包含备份恢复演练，并评估 RTO达成情况。

安全与合规测试覆盖鉴权授权、密钥管理、访问审计、敏感数据脱敏与爬取合规。安全测试包括漏洞扫描、依赖库安全审计、入侵检测（SIEM）联动与敏感信息识别验证。合规测试还要验证审计日志完整性、访问控制策略是否能按权限隔离运维与开发操作，以及对爬取行为的合规审核流程是否生效。

可扩展性与运维测试验证自动伸缩策略、金丝雀发布流程与回滚路径。通过模拟突发流量或资源瓶颈，评估扩容时间、服务恢复时间与告警触发逻辑是否满足 SLA。效能测试用于衡量不同配置对资源消耗与质量指标的影响，帮助制定合理的容量策略与成本预算。
\section{本章小结}
本章围绕Search Agent系统的测试与验证工作展开，从多个维度对系统性能与功能进行系统性评估。首先，在功能测试方面，通过构建典型测试用例，对Prompt解析、实时检索、离线检索及结果融合等核心功能进行逐项验证，确保系统各模块能够按照设计要求稳定运行。在性能测试方面，重点关注系统响应时间、吞吐能力及并发处理能力，通过模拟高负载场景，分析系统在实际应用中的表现，并识别潜在性能瓶颈。同时，在稳定性与可靠性测试中，通过长时间运行测试与异常场景模拟，验证系统在复杂环境下的鲁棒性与容错能力。此外，本章还结合监控系统对关键指标进行实时观测，对系统运行状态进行量化评估。测试结果表明，系统在功能完整性与性能指标方面均达到了预期目标，能够满足实际应用需求。通过本章测试，不仅验证了系统设计与实现的正确性，也为后续优化与实际部署提供了重要依据。